% FILE: 07_open_questions_and_next_work.tex
% PROJECT: crosswalks
% ROLE: type-law-crosswalk-between-standards-and-execution

\section{Open Questions and Next Work}
\label{sec:crosswalks-openq}

\subsection{Known Gaps}
\label{sec:crosswalks-openq-gaps}

Four structural gaps are documented in the crosswalks project evidence record:

\begin{description}
  \item[Gap 1 --- No README:] No \texttt{README.md} exists in \texttt{crosswalks/} (source:
        \texttt{project\_manifest.yaml}, \texttt{readme\_present: false}). The directory is
        navigable only by reading each crosswalk document individually. A README would provide
        a structured index of the eight crosswalks and their projection direction taxonomy.

  \item[Gap 2 --- No explicit receipt:] No \textsc{BLAKE3}-signed receipt chain exists for
        the \texttt{crosswalks/} directory as a unit (\texttt{receipt\_present: false}).
        The \textsc{CROSSWALK\_RECEIPT} cited in the artifact census is ``implicit; per-crosswalk
        dating'' based on file modification timestamps. Under the receipt doctrine defined in
        \texttt{doctrine/RECEIPT\_DOCTRINE.md}, an implicit receipt is not a witnessed receipt.
        The directory cannot claim \textsc{ALIVE} status under the receipt law.

  \item[Gap 3 --- No ontology:] No \textsc{OWL} or Turtle ontology file exists in
        \texttt{crosswalks/} (\texttt{ontology\_present: false}). The crosswalk mappings define
        structural projection laws in prose and table form but do not expose them as a
        machine-queryable ontology. This limits the ability of downstream tools to reason over
        the crosswalk taxonomy programmatically.

  \item[Gap 4 --- No compilation evidence:] The Rust struct definitions in
        \texttt{pm4py\_to\_wasm4pm\_data\_types.md} and \texttt{declare\_to\_wasm\_mapping.md}
        are research specifications. No evidence exists that these structs are present in the
        actual \texttt{wasm4pm-compat} crate source tree. The gap between research specification
        and crate implementation is unverified (source: open questions in evidence extraction).
\end{description}

\subsection{Deferred Gates}
\label{sec:crosswalks-openq-deferred}

Two gates are deferred from the current PARTIAL assessment:

\begin{enumerate}
  \item \textbf{Receipt gate:} Manufacture an explicit \textsc{BLAKE3}-signed receipt for the
        \texttt{crosswalks/} directory. The receipt must bind the eight document content hashes,
        the sealed date (2026-06-01), and the \textsc{ALIVE\_001} verification code
        \texttt{0x0F\_ALL\_GATES\_PASS} as the authority reference. Until this receipt exists,
        the directory cannot be cited as a receipted artifact in the M\&A diligence package.

  \item \textbf{Implementation gate:} Confirm that the Rust structs specified in the crosswalk
        documents (\texttt{WitnessPrecedence}, \texttt{WitnessResponse},
        \texttt{WitnessCoexistence}, \texttt{WitnessSuccession}, \texttt{OcelEventLog},
        \texttt{XesTrace}, \texttt{ActivityName}, \texttt{Timestamp}, \texttt{ResourceId},
        \texttt{TypedLoopNode<ARITY>}) are present and compilable in the
        \texttt{wasm4pm-compat} crate. This gate requires a \texttt{cargo check} or
        \texttt{cargo test} invocation with output demonstrating the types resolve.
\end{enumerate}

\subsection{Research Risks}
\label{sec:crosswalks-openq-risks}

Three research risks are identified:

\begin{description}
  \item[Risk 1 --- \textsc{OR}-gateway completeness:] The \texttt{BPMN\_TO\_WFNET.md} crosswalk
        identifies the \textsc{OR}-gateway as ``the most significant loss'' in the \textsc{BPMN}
        to WF-net direction. The crosswalk specifies that any approximation is semantically
        incorrect for some \textsc{BPMN} processes. If the \texttt{wasm4pm-compat} crate
        implements an \textsc{OR}-gateway approximation rather than a clean refusal, the
        \texttt{BpmnToWfNetRefusal::InclusiveGatewayNotRepresentable} type would be unreachable
        code --- and the crosswalk law would be vacuously satisfied without being enforced.

  \item[Risk 2 --- Lattice join completeness:] The \texttt{bpmn\_to\_declare\_conformance.md}
        and \texttt{declare\_to\_wasm\_mapping.md} crosswalks define lattice join operators for
        four constraint types. The distributed process mining scenario that motivates the join
        (merging independent witness states) is not demonstrated by any test fixture. If the join
        semantics are incorrect for the \textsc{XOR} not-coexistence lattice (where
        $\text{Seen}(A) \sqcup \text{Seen}(B) = \top$), a conformance check over a distributed
        log could produce a false non-violation.

  \item[Risk 3 --- \textsc{ISO} mapping authority:] The \texttt{alive\_001\_to\_iso\_compliance.md}
        crosswalk maps \textsc{ALIVE\_001} gates to four \textsc{ISO/IEC} standards. No formal
        compliance audit or third-party verification of these mappings exists. The mappings are
        author-level research claims, not certified compliance assertions. Citing them as
        regulatory compliance in an M\&A diligence package without qualification is a research
        integrity risk.
\end{description}

\subsection{Next Receipted Work}
\label{sec:crosswalks-openq-next}

The next receipted work for the crosswalks project, in order of priority:

\begin{enumerate}
  \item Manufacture an explicit \textsc{BLAKE3}-signed receipt for the \texttt{crosswalks/}
        directory, binding all eight document content hashes, the sealed date, and the
        \textsc{ALIVE\_001} authority reference. File as
        \texttt{crosswalks/CROSSWALK\_RECEIPT.json}.

  \item Verify implementation of the specified Rust structs in \texttt{wasm4pm-compat} by
        running \texttt{cargo check} against the crate and confirming the eight named types
        resolve without error. Document the output as a compilation witness.

  \item Author a \texttt{crosswalks/README.md} that indexes the eight crosswalks by projection
        direction, loss classification (lossless / conditionally lossless / inherently lossy),
        and \textsc{LossPolicy} variant, providing a navigable taxonomy for downstream
        implementers.

  \item \textbf{[FUTURE WORK]:} Author a Turtle or \textsc{OWL} ontology for the crosswalk
        taxonomy, enabling \textsc{SPARQL} queries over the projection direction graph and
        loss-type lattice. This would make the crosswalk laws machine-queryable and support
        automated conformance of the crosswalk taxonomy itself against the standards it maps.
\end{enumerate}
